\documentclass[12pt,a4paper]{article}

\usepackage{fontawesome5}
\providecommand{\faPencilAlt}{\faIcon{pencil-alt}}

\usepackage{styles/couleurs}
\usepackage{styles/mise_en_page}
\usepackage{styles/boites}
\usepackage{styles/maths}
\usepackage{styles/figures}

\usepackage{enumitem}
\usepackage{amsmath}
\usepackage{tikz}

\newcommand{\lignepointillee}{%
  \noindent\dotfill\par
  \vspace{0.45em}
}

\newcommand{\lignesreponse}[1]{%
  \begin{tcolorbox}[
    enhanced,
    colback=white,
    colframe=NeutreC!40,
    boxrule=0.5pt,
    arc=1.5mm,
    left=3mm,
    right=3mm,
    top=2mm,
    bottom=1mm,
    before skip=2pt,
    after skip=4pt
  ]
    \foreach \i in {1,...,#1}{\lignepointillee}
  \end{tcolorbox}
}

\newcommand{\titreseq}[1]{%
\begin{center}
  {\sffamily\LARGE\bfseries\color{TitreC}
  Séquence 1 -- Statistiques}\\[0.2em]
  {\sffamily\large\bfseries\color{SavoirC}#1}
\end{center}
}

\newcommand{\reponseeleve}[2]{%
\begin{tcolorbox}[
  enhanced,
  colback=NeutreBg,
  colframe=NeutreC!55,
  boxrule=0.6pt,
  arc=2mm,
  left=3mm,
  right=3mm,
  top=2mm,
  bottom=2mm,
  before skip=4pt,
  after skip=4pt
]
{\sffamily\bfseries\color{TitreC}#1}

#2
\end{tcolorbox}
}

\begin{document}

\titreseq{Activité 1.5 -- Peut-on valider ces médianes ?}

\consigne{Pour chaque réponse, décide si tu la valides.
Si tu ne la valides pas, corrige-la et explique l'erreur.}

\reponseeleve{Réponse d'Inès}{
La série étudiée est :
\[
8\ ;\ 4\ ;\ 11\ ;\ 6\ ;\ 7\ ;\ 5\ ;\ 9
\]

Inès écrit :

\begin{quote}
« Il y a 7 valeurs. La valeur du milieu est la 4\ieme{} :
la médiane est donc 6. »
\end{quote}
}

\lignesreponse{3}

\reponseeleve{Réponse de Sami}{
La série ordonnée est :
\[
4\ ;\ 6\ ;\ 7\ ;\ 9\ ;\ 12\ ;\ 14
\]

Sami écrit :

\begin{quote}
« Il y a 6 valeurs. Je prends une des deux valeurs du milieu :
la médiane est 7. »
\end{quote}
}

\lignesreponse{3}

\reponseeleve{Réponse de Zoé}{
La série ordonnée est :
\[
3\ ;\ 5\ ;\ 6\ ;\ 10\ ;\ 12\ ;\ 15
\]

Zoé calcule :
\[
(6+10)\div2=8
\]

puis écrit :

\begin{quote}
« Mais 8 n'apparaît pas dans la série, donc ça ne peut pas être la médiane. »
\end{quote}
}

\lignesreponse{3}

\begin{tcolorbox}[
  enhanced,
  colback=SavoirBg,
  colframe=SavoirC!65,
  boxrule=0.7pt,
  arc=2mm,
  left=3mm,
  right=3mm,
  top=2mm,
  bottom=2mm,
  before skip=6pt
]
{\sffamily\bfseries\color{SavoirC}Pour finir}

Dans la dernière série, comment peut-on vérifier que \(8\) joue bien le rôle de médiane ?
\end{tcolorbox}

\lignesreponse{3}

% Relances enseignant (commentées)
% Inès : « Qu'est-ce qui te gêne ou non dans son raisonnement ? »
% puis « La position d'une valeur a-t-elle un sens dans cette présentation ? »
% puis « Essaie de ranger les valeurs. »
%
% Sami : « Que se passe-t-il au centre quand l'effectif est pair ? »
% puis « Quelles sont les deux valeurs centrales ? »
%
% Zoé : « Une médiane doit-elle forcément appartenir à la série ? »
% puis « Combien de valeurs sont situées de chaque côté de 8 ? »
%
% Réponses attendues :
% Inès : 4 ; 5 ; 6 ; 7 ; 8 ; 9 ; 11, donc médiane 7.
% Sami : (7+9)/2 = 8.
% Zoé : 8 est bien une médiane.

\end{document}
